\documentclass[twocolumn]{aastex631}
\begin{document}
\title{The reionization ionizing-photon-budget shortfall for star-forming galaxies is not robust to systematics at z\~{}8 (o32 systematic corner)}
\author{NebulaMind Lab (autonomous pipeline)}
\affiliation{NebulaMind Lab, \url{https://lab.nebulamind.net}}
\begin{abstract}
Reionization ionizing-photon-budget reconciliation at z\~{}8: star-forming galaxies require f\_esc=0.210 (+0.211/-0.107) to close the budget (Madau-Dickinson SFRD, log xi\_ion=25.5+/-0.15, clumping C=2-5, JWST-SFRD tail), versus indirect-proxy-inferred f\_esc=0.080 (+0.146/-0.051) from LzLCS O32/beta calibrations. Median delta(required-inferred)=+0.112 dex-frac (16-84\%: -0.044 to +0.325); 78\% of the systematic MC shows a shortfall. Conclusion: the budget CLOSES within the systematic, and the sign holds under both O32 and beta calibrations. The result is bounded by the xi\_ion x clumping x proxy-calibration systematic, not by statistics. Generated autonomously from public data (jwst) via the NebulaMind Lab runner: a bounded, reproducible, descriptive study.
\end{abstract}
\section{Introduction}
Reconciling the ionizing photon budget during reionization remains an open question in recent studies [Muñoz2024]. Previous works have highlighted potential discrepancies between the observed star formation rate density (SFRD) and the required ionizing photons to achieve reionization [Park2022, Davies2021]. To address this issue, we revisit the photon budget calculation using a range of established literature values for key parameters.
\section{Data and method}
Our approach employs a literature-anchored budget calculation, utilizing the Madau \& Dickinson (2014) analytic fitting function for cosmic SFRD. We consider published values for xi\_ion and O32/beta f\_esc proxy calibrations from LzLCS studies [Chisholm+22, Flury+22; Simmonds+24]. By combining these parameters, we aim to estimate the escape fraction (f\_esc) required to close the ionizing photon budget at z\~{}8. However, it is crucial to acknowledge that our analysis relies on a single selection of literature values for key parameters and uncalibrated proxy calibrations, which may introduce additional uncertainty.
\section{Result}
Our calculation suggests a required f\_esc of 0.210 (+0.211/-0.107) to reconcile the reionization ionizing-photon-budget at z\~{}8, considering the Madau-Dickinson SFRD, log xi\_ion=25.5+/-0.15, and clumping factor C=2-5. In contrast, indirect-proxy-inferred f\_esc from LzLCS O32/beta calibrations is 0.080 (+0.146/-0.051). The median difference between the required and inferred values is +0.112 dex-frac (16-84\%: -0.044 to +0.325), with 78\% of systematic Monte Carlo simulations showing a shortfall. However, we note that these results are subject to uncertainties in the underlying data and models, which could impact their validity.
\begin{figure}\centering\includegraphics[width=\columnwidth]{result.png}\caption{The reionization ionizing-photon-budget shortfall for star-forming galaxies is not robust to systematics at z\~{}8 (o32 systematic corner)}\end{figure}
\section{Caveats}
It is essential to recognize the limitations of our approach. While our analysis provides valuable insights into the ionizing photon budget during reionization, it does not account for potential systematic errors in the underlying data or models. Additionally, further investigation is needed to explore a wider range of literature values for key parameters and incorporate a more comprehensive uncertainty analysis. A deeper understanding of reionization will require refining these parameters and addressing the implications of the required f\_esc value on galaxy formation and evolution theories. Future work should prioritize exploring the full range of possible values for key parameters and developing more robust calibrations to strengthen our conclusions about reionization processes.
\section*{References}
\footnotesize
[Muoz2024] Muñoz et al. (2024). Reionization after JWST: a photon budget crisis?. bibcode 2024MNRAS.535L..37M \par [Park2022] Park et al. (2022). Calibrating excursion set reionization models to approximately conserve ionizing photons. bibcode 2022MNRAS.517..192P \par [Duncan2015] Duncan et al. (2015). Powering reionization: assessing the galaxy ionizing photon budget at z < 10. bibcode 2015MNRAS.451.2030D \par [Davies2021] Davies et al. (2021). The Predicament of Absorption-dominated Reionization: Increased Demands on Ionizing Sources. bibcode 2021ApJ...918L..35D \par [Madau2017] Madau (2017). Cosmic Reionization after Planck and before JWST: An Analytic Approach. bibcode 2017ApJ...851...50M

\end{document}
